\chapter{Introduction}
\label{ch:introduction}

\section{Context and Motivation}
\label{sec:intro-context}
Agentic AI systems are now used daily to write and maintain production software.
Developers, are still understanding how to use them well. 
Software construction is not a single activity but a sequence of distinct phases
(planning, implementation, debugging, review, and maintenance) each with
its own goals, its own measure of correctness, and its own failure modes. A
general-purpose language model that performs well in one phase offers no
guarantee of doing so in another, which makes phase-specific study of agent
behaviour necessary rather than optional.

This thesis concerns the debugging phase. Debugging is expensive: before a single
line of a fix is written, a developer typically writes or reads a bug report,
gathers diagnostic information, construct or not, test cases that exposes the failure,
and investigates the root cause. That workflow is repetitive and time-consuming in development cycles, which makes it an
attractive target for automation, and a demanding one.

Recent advances in large language models (LLMs) and agent-based software
engineering systems have shown promise on tasks such as code generation, test
generation, and program repair. Real-world bug fixing nonetheless remains
difficult. On SWE-bench, a benchmark of 2,294 genuine GitHub issues introduced by
Jimenez et al.~\cite{Jimenez2023SWEbench}, the SWE-agent system of Yang et
al.~\cite{Yang2024SWEagent} resolves approximately 12.5\% of issues. A
state-of-the-art result at the time of publication, and still a reminder of how
far automated repair is from reliability on real code.

One finding from recent work is particularly relevant to that gap: the presence
of a \emph{Bug Reproduction Test} (BRT, a test that fails on the buggy
revision and passes once the bug is fixed) materially improves repair
outcomes. In a study on Google's internal codebase, Cheng et
al.~\cite{Cheng2025Agentic} showed that supplying an automated program repair
system with high-quality BRTs increased the proportion of bugs reaching a
plausible fix.

That finding has not, however, been carried into agent \emph{architecture}.
Contemporary repair agents remain reactive: they attempt a patch first and
validate afterwards, so the test, when it appears at all, arrives as
after-the-fact confirmation rather than as a constraint on the search. Also they have only testes on internal databse and issues, and not on real bug issues as SWE-Bench lite provides. This
thesis changes the order, requiring an agent to produce and validate a bug
reproduction test using the bug report and repository context \emph{before} any
fix is attempted. Mirroring test-driven practice in human development, and
using the agent's own code comprehension to sharpen the objective it is then
asked to try to come up with a soloution for the issue.

The theory is that a reproduction test functions as an executable
specification of the failure, constraining the space of acceptable patches. By
forcing reproduction first, the fix stage receives an earlier and sharper
feedback signal than post-hoc validation can provide, which should reduce
unfocused edits and improve repair precision.

\section{Problem Statement}
\label{sec:intro-problem}
Current agents do not have a standard strcuture to solve bugs, it will depend on user prompt, without holding a lot of contrainsts on the failure. That is the gap we working on. How do we engineer agentic systems that fix more bugs by adopting a test-first order rather than reactive validation? The thing with reactive agent is that they have no exectable specifications, the agent relies on its reading of the issue, which can lead to misinterpratation, unfocused edits, regressions in unrelated code, and a self-declared stopping criterion. The TDD comes to guide the LLM on throught those issues, so it can always test if what it's being changed its being done on the right place and not breaking other parts of the code. Test generation from LLM is well studied, however it was not done all in one unique pipeline, rather are different systems that work separatelly. We focus on how making TDD mandatory for bug fixing in one agentic system only can make a difference.

% \section{Thesis Statement and Approach}
% \label{sec:intro-thesis}

\section{Research Questions}
\label{sec:intro-rqs}

The literature brings test-generation quality, repair capability, and
orchestration each affect how many real-world issues an agentic system
resolves~\cite{Meng2024Empirical}. It does not establish how those components
should be composed, nor in which order a practitioner should run them. The three
questions below move from the overall effect of ordering, to the properties of
the test artifact through which any such effect must be carried, to the pipeline
characteristics that sustain it under a fixed budget.

\paragraph{RQ1 --- Does a test-first agentic pipeline, in which repair is
withheld until a generated test has been observed to fail on the buggy revision,
resolve more real-world issues than a matched reactive pipeline?}
This question is answered by a paired comparison on SWE-bench Lite in which every
sampled instance is processed by both arms under the same model, tool set,
iteration and token budgets, container, and scoring harness; the arms differ only
in the presence of the reproduction phase and of the tool that executes its
output. \emph{Reactive} denotes the prevailing design, exemplified by
SWE-agent~\cite{Yang2024SWEagent}, in which an agent may run tests but is under
no obligation to establish a failing one before editing begins. The primary
outcome is resolved-rate: the benchmark's held-out failing tests pass and its
previously passing tests are preserved.

\paragraph{RQ2 --- Which pipeline and orchestration characteristics improve fix
quality and reliability?}
An agentic pipeline admits many designs: how many agents participate, which tools each may call, how they pass work to one another, and what evidence is required before a test or a fix is accepted. This question is answered by ablating the test-first pipeline one component at a time, holding the model, the instance sample, the container, and the scoring harness fixed: the criterion by which a generated test is accepted without access to the reference patch, the per-instance attempt budget, agent self-review of the test and of the fix, post-fix adversarial verification, and best-of-$N$ test selection. Each variant is reported as its change in resolved-rate against its measured cost per instance. Two ablations produced negative results; they are reported as findings rather than omitted. 

\paragraph{RQ3 --- Can open-weight models running on local hardware sustain the test-first pipeline, and at which stage does the trajectory break down?}
Open-weight models in the 7–14B range are attractive for repair work: they are free to run, and the code never leaves the machine. Whether they can act as agents is a separate question from whether they can write code. This question is answered by running the same pipeline, prompts, tools, and instances as the frontier arm, with the context window and iteration budget held fixed and reported, and by recording where each trajectory ends: whether the agent reaches a terminal submission within budget, whether the test it submits reproduces the failure on the buggy revision, and whether the fix it then produces resolves the issue. Reporting the stage at which the pipeline stops, rather than resolved-rate alone, distinguishes a scaffolding limit from a capability limit — a distinction that determines whether the gap is closable by prompt design. 


